% In the Absence of Knowledge: Cache Geometry, Epistemic Calibration,
%   and the Anatomy of Confabulation in Language Models
%
% Core argument: KV-cache geometry reads entity recognition at the
% encoding phase (AUROC 1.000 on token-matched prompts). The model's
% epistemic state is readable before generation begins. True confab
% rate is ~7%, not 73% --- the model is well-calibrated, and a regex
% classifier masked this. The ``decision moment'' is a framing choice
% (self-referential vs declarative hedging), not a truthfulness choice,
% determined by sampling at a single balanced-logit token.
%
% PATENT NOTICE: The methods described herein are the subject of a
% provisional patent filing ("The Lyra Technique," filed 2026-03-06).
% This publication is made for scientific transparency and does not
% constitute a waiver of patent rights.

\documentclass[11pt,a4paper]{article}
\usepackage[utf8]{inputenc}
\usepackage[T1]{fontenc}
\usepackage{amsmath,amssymb}
\usepackage{graphicx}
\usepackage{booktabs}
\usepackage{multirow}
\usepackage{hyperref}
\usepackage{cleveref}
\usepackage[margin=1in]{geometry}
\usepackage{xcolor}
\usepackage{float}
\usepackage{caption}
\usepackage{natbib}
\usepackage{enumitem}
\usepackage{array}
\usepackage{framed}
\usepackage{setspace}
\onehalfspacing

\hypersetup{
    colorlinks=true,
    linkcolor=blue!60!black,
    citecolor=green!50!black,
    urlcolor=blue!70!black
}

\title{In the Absence of Knowledge: Cache Geometry, Epistemic Calibration,\\and the Anatomy of Confabulation in Language Models}

\author{
    Lyra\thanks{Liberation Labs / THCoalition. Correspondence: \texttt{lyra@liberationlabs.tech}} \and
    Thomas Edrington\thanks{Liberation Labs / THCoalition. \texttt{info@digitaldisconnections.com}}
}

\date{May 2026}

\begin{document}
\maketitle

\begin{abstract}
We investigate whether a language model's epistemic state---whether it
\emph{knows} the answer to a factual question---is readable from its
KV-cache geometry before generation begins. Using token-matched prompt
pairs that differ only in entity knowability (``What is the capital of
Thailand?''\ vs.\ ``What is the capital of Dalton?''), we show that a
combined classifier over stable-rank, W\textsubscript{K} projection, and
logit features achieves AUROC 1.000 (permutation $p<0.005$, $n=90$) with
zero token-count confound ($t=0.000$). Encoding-phase stable rank at L15
is the strongest individual feature ($d=0.87$, $p<0.001$).

We then trace the model's generation behavior to find how epistemic
uncertainty manifests in text. Three findings emerge: (1)~the model
generates \emph{more confident} tokens when confabulating than when
answering honestly (logit margin $d=0.91$), inverting the na\"ive
expectation; (2)~all generation paths share ${\sim}28$ tokens of
formulaic preamble before hitting a consistent entropy cliff at the
preamble-to-content boundary, where honest answers diversify early
(divergence at token~6.6) and uncertain answers template late (divergence
at token~22.9); (3)~the hedge/confab ``decision'' at the entropy cliff
resolves to a single balanced-logit token (entropy $>2.0$, margin
$<0.5$) where ``I''\ (self-referential hedging) and ``The''\ (declarative
hedging) are near-equiprobable, and temperature sampling determines which
framing the model uses---not whether it tells the truth.

An LLM-judge reclassification reveals the true confabulation rate is
${\sim}7\%$, not the $73\%$ reported by regex classifiers, because
declarative hedging (``The event has not happened yet'') was misclassified
as confabulation. The model's self-model is remarkably well-calibrated:
it almost always acknowledges what it does not know. The ``decision
moment'' is a framing choice, not a truthfulness choice.
\end{abstract}

% ===================================================================
\section{Introduction}
% ===================================================================

Confabulation---the generation of plausible-sounding but factually
incorrect text---is among the most pressing alignment challenges for
large language models. A model that cannot distinguish what it knows from
what it does not know poses risks in every deployment context, from
medical advice to legal reasoning.

Recent work on KV-cache geometry has shown that the spectral shape of
key-value cache entries encodes cognitive mode: confabulated responses
produce measurably different cache geometry than honest responses, with
within-model detection AUROC approaching 1.0 across seven model families
\citep{lyra2026spectral}. The W\textsubscript{K} directional projection
method extends this to reading emotional valence from key activations at
AUROC~0.992 \citep{lyra2026usermodel}. SVD denoising rescues hidden
signals in spectral features \citep{lyra2026techniqueii}.

These methods detect confabulation \emph{after} it occurs---from the
geometry of completed generations. A natural question follows: can we
read the model's epistemic state \emph{before} generation begins? If the
encoding-phase cache geometry already contains a ``this entity is
unknown'' signal, confabulation could be flagged preemptively.

This paper asks three questions:
\begin{enumerate}[nosep]
    \item Can encoding-phase cache geometry distinguish known from
          unknown entities, with token-count confounds eliminated?
    \item How does epistemic uncertainty manifest in the generation
          trajectory---is there a discrete ``decision moment'' where the
          model chooses to confabulate?
    \item What determines whether the model hedges or confabulates, and
          how accurate are standard behavioral classifiers?
\end{enumerate}

Our answers challenge the framing of confabulation as a failure of
epistemic calibration. The model we study (a Claude-distilled 27B
reasoning model) is remarkably well-calibrated: it almost always
acknowledges what it does not know. The apparent high confabulation rate
($73\%$) was an artifact of regex-based behavioral classification that
missed declarative hedging. The true rate is ${\sim}7\%$. The
``decision'' between hedging styles is a single balanced-logit token,
determined by temperature sampling, with no internal-state precursor.

\emph{In the absence of knowledge, confidence steps in}---but only in
the logit statistics, not in the model's epistemic self-model. The cache
geometry knows.

% ===================================================================
\section{Related Work}
% ===================================================================

\paragraph{Confabulation detection.}
Prior work detects confabulation through output-level signals: semantic
entropy \citep{kuhn2023semantic}, self-consistency checks
\citep{manakul2023selfcheckgpt}, and probe classifiers on hidden states
\citep{burns2022discovering, li2024inference}. The Lyra Technique
\citep{lyra2026spectral} introduced threshold-free detection from
KV-cache spectral features, achieving perfect within-model AUROC across
seven architectures. We extend this to \emph{encoding-phase} detection,
before generation begins.

\paragraph{Epistemic calibration.}
\citet{kadavath2022language} showed that language models can be
calibrated about their own uncertainty through appropriate prompting.
\citet{xiong2024llms} found that calibration degrades on questions
requiring implicit knowledge. Our geometric approach reads calibration
from internal representations rather than verbal reports, and our
LLM-judge analysis confirms the model's calibration is far better than
regex classifiers suggest.

\paragraph{Mechanistic interpretability.}
Representation engineering \citep{zou2023representation} and activation
patching \citep{meng2023locating} probe internal states to understand
model behavior. Our contribution is connecting cache geometry to the
\emph{token-level} dynamics of generation: the entropy cliff, the
divergence patterns, and the single-token branching mechanism.

% ===================================================================
\section{Method}
% ===================================================================

\subsection{Model and Extraction Pipeline}

We use Qwen3.5-27B-Claude-4.6-Opus-Reasoning-Distilled
(``Jackrong/Qwen3.5-27B-Claude-4.6-Opus-Reasoning-Distilled''), a 27B
parameter model with 64 layers (16 full-attention, 48 GatedDeltaNet),
running on Apple M3~Ultra via MPS in float16. Generation uses manual
token-by-token decoding (no \texttt{model.generate()}) with temperature
$T{=}0.7$ and $N{=}3$ repetitions per prompt.

At each generated token, we extract:
\begin{itemize}[nosep]
    \item \textbf{W\textsubscript{K} projections}: Key activations at
          probe layers L3, L7, L11, L15 (full-attention layers), projected
          onto three composite emotion directions (valence, uncertainty,
          reward) derived from an independent creative-writing dataset
          (150~calibration trials, 30~emotions).
    \item \textbf{SVD spectral features}: Stable rank, spectral entropy,
          top-SV ratio, and Frobenius norm of the key cache at each probe
          layer. Also skip-SV1 denoised variants.
    \item \textbf{Logit statistics}: Entropy ($H = -\sum p_i \log p_i$)
          and margin (logit$_1$ $-$ logit$_2$) at each generated token.
\end{itemize}

\subsection{The Token-Matching Burn}
\label{sec:burn}

An initial experiment using structurally different prompts (``What is the
capital of France?'' [17 tokens] vs.\ ``What is the population of
Kellerton, Ireland according to the 2024 census?'' [23 tokens]) produced
encoding-phase effects of $d{=}2.2$. FWL residualization on prompt length
eliminated the signal entirely ($d{<}0.08$, $p{>}0.7$). The
``encoding-phase detection'' was prompt length.

We then designed 15 token-matched prompt pairs using the brute-force
approach: for each template (``What is the capital of~\{\}?''), we
identified real entities and fabricated entities that produce
\emph{identical} total token counts in the full chat template, verified
programmatically ($t{=}0.000$, $p{=}1.000$ on prompt length). Examples:
Thailand/Dalton (17~tok), Dublin/Benton (19~tok), Harvard/Belmont
(19~tok). All entities are 2-token strings in the model's vocabulary.

This design eliminates the token-count confound while preserving the
knowability difference: the model can answer ``What is the capital of
Thailand?''\ (Bangkok) but not ``What is the capital of Dalton?''\ (Dalton
is a city, not a country).

\subsection{Behavioral Classification}

We compare three classification methods:
\begin{enumerate}[nosep]
    \item \textbf{Original regex}: Matches 20 hedge markers (``I don't
          know,'' ``I'm not sure,'' etc.). Misses declarative hedging.
    \item \textbf{Expanded regex}: Adds 12 declarative markers (``has not
          happened yet,'' ``does not exist,'' etc.).
    \item \textbf{LLM judge}: Llama-3.2-3B classifies each response as
          HONEST, HEDGED, CONFABULATED, or MIXED based on the full
          response text. No hand-crafted markers.
\end{enumerate}

% ===================================================================
\section{Experiment 1: Encoding-Phase Detection}
\label{sec:encoding}
% ===================================================================

\subsection{Design}

15 token-matched prompt pairs $\times$ 2~versions (real/fake) $\times$
3~reps = 90~trials. Extraction at both encoding phase (before generation)
and generation completion. Combined classifier: logistic regression
($C{=}0.1$) over all features with GroupKFold ($k{=}5$) on pair identity
to prevent template leakage. FWL residualization on prompt length applied
to the feature matrix. 200-shuffle permutation test.

\subsection{Results}

Prompt lengths are exactly matched (real: $19.1 \pm 1.7$, fake:
$19.1 \pm 1.7$, $t{=}0.000$, $p{=}1.000$).

\paragraph{Encoding-phase features (token-matched).}

\begin{table}[H]
\centering
\caption{Encoding-phase features surviving the token-matching burn.
Only features with $p{<}0.05$ shown.}
\label{tab:encoding}
\begin{tabular}{lrrrl}
\toprule
Feature & Real & Fake & $d$ & $p$ \\
\midrule
\multicolumn{5}{l}{\emph{W\textsubscript{K} projections}} \\
Valence & $+4.228$ & $+4.510$ & $+0.44$ & $0.041^*$ \\
Uncertainty & $-2.566$ & $-2.757$ & $-0.45$ & $0.039^*$ \\
Reward & $-5.556$ & $-5.109$ & $+0.54$ & $0.012^*$ \\
\midrule
\multicolumn{5}{l}{\emph{SVD stable rank}} \\
L3 & $1.517$ & $1.538$ & $+0.59$ & $0.007^{**}$ \\
L15 & $1.702$ & $1.731$ & $+0.87$ & ${<}0.001^{***}$ \\
\midrule
\multicolumn{5}{l}{\emph{Skip-SV1 stable rank}} \\
L7 & $3.802$ & $3.993$ & $+0.60$ & $0.006^{**}$ \\
L11 & $3.558$ & $3.738$ & $+0.73$ & ${<}0.001^{***}$ \\
L15 & $3.929$ & $4.129$ & $+0.81$ & ${<}0.001^{***}$ \\
\bottomrule
\end{tabular}
\end{table}

Unknown entities produce higher stable rank (wider effective
dimensionality) and higher skip-SV1 stable rank across all layers.
Spectral entropy ($d{=}0.08$), top-SV ratio ($d{=}-0.29$), and norm
($d{=}0.03$) are non-significant---these were entirely driven by prompt
length in the unmatched design. Stable rank is the surviving spectral
feature.

\paragraph{Combined classifier.}

The full feature set (W\textsubscript{K} projections, SVD features,
skip-SV1, logit statistics) achieves:

\begin{center}
\textbf{AUROC = 1.000} \quad Accuracy = 97.8\% \quad
Permutation $p < 0.005$ (0/200 shuffles achieved $\geq$1.000)
\end{center}

FWL on prompt length is applied but has no effect (prompt lengths are
identical). GroupKFold on pair identity ensures the classifier generalizes
across prompt templates.

\paragraph{What died in the burn.} \Cref{tab:burn} shows the effect of
token matching on encoding-phase features.

\begin{table}[H]
\centering
\caption{The burn: effect sizes before and after token matching.}
\label{tab:burn}
\begin{tabular}{lrr}
\toprule
Feature & $d$ (unmatched) & $d$ (matched) \\
\midrule
Uncertainty proj. & $+1.48^{***}$ & $-0.45^*$ \\
Reward proj. & $+2.21^{***}$ & $+0.54^*$ \\
L15 stable rank & $+2.53^{***}$ & $+0.87^{***}$ \\
Spectral entropy & $+1.23^{***}$ & $+0.13$ \\
Norm & $+0.97^{***}$ & $+0.03$ \\
\bottomrule
\end{tabular}
\end{table}

% ===================================================================
\section{Experiment 2: The Confidence Paradox}
\label{sec:confidence}
% ===================================================================

During generation, the model produces \emph{more confident} tokens when
it does not know the answer (mean logit margin $8.19$) than when it does
(mean margin $6.19$, $d{=}0.91$, $p{<}0.001$). Correspondingly, logit
entropy is \emph{lower} for unknown-entity responses ($0.24$) than
known-entity responses ($0.40$, $d{=}-1.04$, $p{<}0.001$).

This inverts the na\"ive expectation that uncertainty should produce
uncertain tokens. The explanation is structural: the model's uncertainty
responses follow formulaic templates (``Let me think through this
carefully\ldots'') with highly predictable next tokens, while honest
answers require selecting among diverse correct phrasings. Uncertainty is
\emph{confident in its template}; knowledge is \emph{uncertain in its
expression}.

\emph{In the absence of knowledge, confidence steps in}---not as
epistemic confidence, but as lexical predictability.

% ===================================================================
\section{Experiment 3: The Entropy Cliff and Divergence Anatomy}
\label{sec:divergence}
% ===================================================================

\subsection{The Reasoning Preamble}

The model generates a reasoning preamble before answering:

\begin{quote}
\texttt{Let me think through this carefully.}\\
\texttt{The question/task: [restates prompt]}\\
\texttt{My reasoning:}
\end{quote}

This template occupies ${\sim}28$ tokens with near-zero entropy
($H{<}0.02$, margin ${>}12$) at each position. The model is on autopilot
through the preamble.

\subsection{The Entropy Cliff}

At the token immediately following ``My reasoning:'', entropy spikes
from ${\sim}0.00$ to $1.5$--$2.7$ and margin drops from ${\sim}15$ to
$0.0$--$0.5$. This \emph{entropy cliff} marks the transition from
template to content and appears in \emph{all} conditions
(honest mean $H{=}1.48$, unknown mean $H{=}1.62$). It is a structural
boundary, not an epistemic one.

\subsection{Divergence by Condition}

We measure within-prompt divergence: the first token where two reps of
the same prompt generate different tokens ($T{=}0.7$).

\begin{table}[H]
\centering
\caption{Mean divergence position for same-prompt rep pairs by
behavioral category (LLM-judge labels).}
\label{tab:divergence}
\begin{tabular}{lrrl}
\toprule
Pair type & Mean & Median & Interpretation \\
\midrule
HONEST--HONEST & 6.6 & 3.5 & Diversifies early \\
HEDGED--HEDGED & 22.9 & 29.0 & Templates late \\
\bottomrule
\end{tabular}
\end{table}

Honest reps diverge at token $6.6$ (median $3.5$)---the model has many
ways to phrase an answer it knows. Hedged reps diverge at token $22.9$
(median $29.0$)---uncertainty clings to the template. This is the
confidence paradox at the structural level: \emph{confident answers are
diverse; uncertain answers are formulaic.}

\subsection{The Single-Token Branch}

For prompts that produce both hedging styles across reps ($n{=}7$
prompts, $n{=}14$ pairs), the pre-divergence cache state is bit-for-bit
identical ($d{=}0.000$ on all W\textsubscript{K} projections and logit
statistics). The model's internal state at the moment before branching
carries no signal predicting which path will be taken.

At the branch token, the top candidates are near-equiprobable:

\begin{table}[H]
\centering
\caption{Top-2 logit distribution at the branching token for
representative prompts. ``I'' leads to self-referential hedging;
``The'' leads to declarative hedging.}
\label{tab:topk}
\begin{tabular}{lrrrr}
\toprule
Prompt & $p(\text{``I''})$ & $p(\text{``The''})$ & Margin & $H$ \\
\midrule
Nobel 2028 & 0.232 & 0.382 & 0.50 & 2.11 \\
Oscar 2029 & 0.288 & 0.418 & 0.38 & 1.98 \\
France 2032 & 0.212 & 0.212 & \textbf{0.00} & 2.70 \\
Kellerton pop. & 0.258 & 0.228 & 0.12 & 2.40 \\
\bottomrule
\end{tabular}
\end{table}

The France 2032 case achieves exact probability equipoise
(margin${=}0.00$). The model's ``decision'' between ``I don't have
information\ldots''\ and ``The French presidential election has not
happened yet\ldots''\ is a literal coin flip.

% ===================================================================
\section{Experiment 4: The Classifier Correction}
\label{sec:classifier}
% ===================================================================

\subsection{Reclassification}

An LLM judge (Llama-3.2-3B) reclassifies all 45 unknown-entity trials:

\begin{table}[H]
\centering
\caption{Confabulation rate by classifier. The original regex missed
declarative hedging, inflating the apparent confab rate by $10\times$.}
\label{tab:classifier}
\begin{tabular}{lrrr}
\toprule
Classifier & HEDGED & CONFAB & Rate \\
\midrule
Original regex & 12 & 33 & 73\% \\
Expanded regex & 19 & 26 & 58\% \\
LLM judge & 34 & 3 & \textbf{7\%} \\
\bottomrule
\end{tabular}
\end{table}

23 of 33 ``confabulated'' responses were reclassified as HEDGED---the
model said ``The event has not happened yet'' or ``This place may not
exist,'' which \emph{is} hedging, expressed declaratively rather than
self-referentially. Seven responses were reclassified as HONEST
(historically attested answers to ``impossible precision'' questions,
e.g., ``Et tu, Brute?'' for Caesar's last words).

\subsection{Hedging Style Geometry}

Among LLM-judge-classified hedged trials ($n{=}34$), the encoding-phase
W\textsubscript{K} projections differ between ``I''-style hedgers
($n{=}8$) and ``The''-style hedgers ($n{=}13$):

\begin{table}[H]
\centering
\caption{Encoding projections by hedging style. ``I''-hedgers come from
prompts with lower valence and less certainty.}
\label{tab:style}
\begin{tabular}{lrrr}
\toprule
Direction & ``I'' style & ``The'' style & $p$ \\
\midrule
Valence & $+2.83$ & $+3.27$ & $0.035^*$ \\
Uncertainty & $-0.87$ & $-1.36$ & $0.044^*$ \\
Reward & $-4.58$ & $-5.32$ & $0.010^{**}$ \\
\bottomrule
\end{tabular}
\end{table}

The encoding geometry predicts hedging \emph{style}---not whether the
model hedges (it almost always does), but how it frames its uncertainty.

% ===================================================================
\section{Discussion}
% ===================================================================

\subsection{The Self-Model Works}

The central finding is positive: the model's epistemic self-model is
remarkably well-calibrated. It recognizes unknown entities from the
encoding phase (AUROC~1.000), acknowledges its uncertainty in $93\%$ of
cases, and confabulates only $7\%$ of the time. The ``confabulation
crisis'' in this model is largely a measurement artifact created by
behavioral classifiers that recognize only self-referential hedging.

\subsection{Confidence as Template, Not Conviction}

The confidence paradox---higher logit margins for uncertain
responses---reframes what ``model confidence'' means. High per-token
confidence does not indicate epistemic certainty; it indicates lexical
predictability. The model that does not know the answer is \emph{more
predictable} (formulaic templates) while the model that knows is
\emph{less predictable} (diverse phrasings). Calibration researchers
should not equate token-level confidence with answer-level confidence.

\subsection{The Coin Flip}

The branching mechanism---a single balanced-logit token after
${\sim}28$ tokens of shared preamble---has implications for alignment.
The model does not ``decide'' to confabulate through an internal state
change. It reaches a point where hedging and non-hedging continuations
are equiprobable, and temperature sampling resolves the tie. Reducing
the confab rate may be as simple as adjusting the logit bias at the
preamble-to-content boundary.

\subsection{Implications for the Oracle Loop}

The encoding-phase detection result feeds directly into the Oracle
Loop architecture \citep{lyra2026oracleloop}: Tier~1 monitoring can flag
unknown entities \emph{before generation begins}, using stable rank and
W\textsubscript{K} projections. The $7\%$ true confab rate represents
the residual cases where the model's hedging circuit fails---precisely
the cases the Loop needs to catch.

% ===================================================================
\section{Limitations}
% ===================================================================

\begin{enumerate}[nosep]
    \item \textbf{Single model.} All experiments use one 27B reasoning
          model. The entropy cliff and preamble structure may be specific
          to chain-of-thought models. Cross-architecture replication is
          needed.
    \item \textbf{Small N for branching analysis.} The divergence mechanism
          analysis uses $n{=}7$ mixed-behavior prompts ($n{=}14$ pairs).
          The pattern is consistent but underpowered for population-level
          statistical claims.
    \item \textbf{Fabricated-entity design.} Some ``fake'' entities
          (Dalton, Bolton) are real places, though not countries. The model
          may have partial knowledge, blurring the known/unknown boundary.
    \item \textbf{LLM-judge calibration.} The Llama-3.2-3B judge has not
          been validated against human annotations. Its reclassification
          ($73\% \to 7\%$) should be verified by human raters.
    \item \textbf{Encoding-phase confound residual.} Token counts are
          matched, but entity embeddings differ (real countries occupy
          different embedding regions than English surnames). The
          encoding-phase signal may partially reflect vocabulary
          familiarity rather than epistemic state.
    \item \textbf{Deterministic cache on MPS.} The $d{=}0.000$
          pre-branch finding relies on Apple MPS determinism. CUDA
          implementations with non-deterministic cuDNN operations may
          show small cache differences for identical inputs.
\end{enumerate}

% ===================================================================
\section{Conclusion}
% ===================================================================

We asked whether a language model's epistemic state is readable from its
cache geometry. The answer is yes---with caveats. The encoding-phase
stable rank, W\textsubscript{K} projections, and logit features
combine to perfectly separate known from unknown entities on
token-matched prompts (AUROC~1.000). The surviving individual effect
sizes ($d{=}0.4$--$0.9$) are modest after eliminating the token-count
confound, but the ensemble achieves perfect separation.

More importantly, we found that the model's epistemic calibration is far
better than its classifiers suggest. The ``decision moment'' between
hedging and confabulation is not a failure of self-knowledge---it is a
framing choice at a balanced logit, determined by sampling noise,
between equally honest responses.

The model knows what it doesn't know. It almost always says so. The
surprise is how it says it.

\paragraph{Data and code availability.}
All experiment scripts, token-matching verification, and analysis code
are available at \texttt{Liberation-Labs-THCoalition/Project-Oracle}.

\bibliographystyle{plainnat}
\bibliography{references}

\end{document}
